We have built and executed an end-to-end, published-value-anchored forward
model of the reconstructed-energy background spectra for a thin single-sided
germanium QPD wafer. Starting from the deposited-energy spectra of the three
dominant surface channels---reactor CEvNS, through-going cosmic muons, and the
radiogenic environmental-gamma continuum---we placed all three on a single
unified phonon energy scale with no ionization quenching, and folded them
through a common bandwidth-limited response matrix $R(E_{\rm rec}\,|\,E_{\rm dep})$
for two quasiparticle-trapping designs, Ta$\to$Al and Al$\to$Hf. The output is
the reconstructed-energy differential rate $dR/dE_{\rm rec}$ that a $25$~kHz
readout actually records, rather than the deposited energy that is physically
inaccessible to it.

Three findings organize the result. The reactor-CEvNS signal reconstructs
\emph{linearly} to tens of eV, peaking near $42$~eV in $E_{\rm rec}$ with
$\sim$$85\%$ of its counts below $100$~eV: this is the flagship signal band, and
it lies below the sub-$100$~eV thresholds now being probed by reactor-CEvNS
experiments~\cite{conus2025}. Cosmic muons, depositing $\sim$$1.5$--$197$~MeV,
drive the readout deep into saturation and collapse onto a single tens-of-keV
pile-up ($\approx$$18.8$~keV for Ta$\to$Al, $\approx$$15.0$~keV for Al$\to$Hf)---an
instrumental artifact of the modeled ceiling, not a physical line. The
environmental-gamma Compton continuum fills the region between the two and,
overlapping the signal band directly, is the dominant reducible background. The
raw event rate, though dominated by the two saturating channels, yields a
pileup occupancy of only $R\tau\approx3.3\times10^{-5}$, so quiescent
reconstruction is not precluded. Every channel is anchored: the deposited CEvNS
rate reproduces the Billard germanium benchmark to within
$2.4\%$~\cite{billard2017}, the muon flux matches the Particle Data Group
sea-level value within its $\sim$$30\%$ spread, and the Compton edges are fixed
exactly by scattering kinematics.

We state the scope plainly. This is a stage-one forward-model feasibility
study. It establishes the reconstructed-energy regime in which a QPD germanium
wafer would operate; it does \emph{not} compute a sensitivity and makes no claim
of detectability or discovery significance. Four modeling uncertainties bound
the interpretation and, deliberately, three of them act only on the saturated
backgrounds while the flagship result---that reactor CEvNS reconstructs to tens
of eV---remains the most robust feature: the unbenchmarked saturated-regime
response shape, the factor-of-$\sim$2 site-dependent gamma normalization, the
cited-but-not-reproduced sub-$1.8$~MeV reactor flux, and the crossover deposit
energy carried as a $\sim$$50$~eV--$13$~keV band.

Each caveat maps to a concrete next step. A direct QPD saturation calibration
would replace the modeled ceiling shape; an on-site gamma survey would fix the
absolute Compton normalization; a thin-wafer sensor-sharing measurement would
turn the crossover band into a line; and a digitized low-energy reactor flux
would collapse the signal-side band. Combined with an exposure model and a
noise-and-resolution treatment of the reconstruction floor, these define the
path from the present feasibility map to a quantitative QPD reactor-CEvNS
sensitivity, extending the phonon-scale reach of the QPD concept~\cite{ramanathan2026}
into the reactor-antineutrino program.
